\begin{helper}[Step 1 zero-divider PE branch]
Let $T$ be a positive-rung ETT of rung $n$ on a finite loopless
multigraph with individual edge identities, generating coloring $\psi$,
maximum property, properness, and exactly the distinguished edge $e_0$
uncolored.  Assume uniformly for every positive-rung maximum ETT of rung
strictly below $n$ the complete six assertions used in Step 1:
terminal elementarity; stable interchangeability of the final ladder
closure, including isolated path components; PE stability descent unconditionally before each support and, when the
first-appearance index is greater than one, to the preceding ladder stage;
stored pairs at distinct supports with the same first appearance are disjoint; the terminal PE component intersection assertion; the
defective-endpoint precedence assertion; and the terminal mod-coloring
assertion.  Assume that the terminal extension kind is PE.

Let $B$ be a hierarchy base and $H$ a hierarchy over $B$ with zero
dividers.  Put $A=H.\mathit{starred}(0)$, the exact hierarchy core.  Let
$a\ne b$ be colors such that $a$ is missing and singleton-closed on $A$.
Suppose there are exact $(A,\psi,\{a,b\})$-exit paths $P_1,P_2$ with
distinct outside endpoints and disjoint incidence-walk supports.  Then
these data are contradictory.
\end{helper}
\begin{proof}
Write $n$ for the rung number, $T_n$ for the tree in series slot $n$,
$\psi_{n-1}$ for the coloring in slot $n-1$, and $D_i$ for the exact
computed stored-color set in slot $i$.  Thus the predecessor uses
$D_{n-1}$, while the terminal PE assertions use $D_n$.  The
generating-coloring field of \noderef{ExtendedTashkinovTree} identifies
the displayed coloring $\psi$ with series color slot $n$.  Since
the divider count is zero, the value of
\noderef{HierarchyStarredStage} at its only index is the core stored in
the given \noderef{HierarchyBase}.  Thus the tree called $A$ in the
statement is exactly that core.

We first record two consequences of each alleged exit path.  Unfold
\noderef{TreeExitPath}.  Its node list starts at the outside vertex-node
$a_i$, ends at the tree vertex-node $b_i$, has length at least two, is
simple, and contains no other vertex-node of $A$.  By the bipartite
adjacency definition in \noderef{TwoColorIncidenceGraph}, the node just
before $b_i$ is an exact edge identity $f_i$ colored $a$ or $b$ and
incident with $b_i$.  The vertex-node before $f_i$ is outside $A$,
whereas $b_i$ is inside.  Hence the endpoint description in
\noderef{BoundaryEdges} gives $f_i\in\partial_G(A)$, and
\noderef{BoundaryColorEdges} gives
\[
 f_i\in\partial_{\psi,a}(A)\quad\hbox{or}\quad
 f_i\in\partial_{\psi,b}(A).                         \tag{1}
\]
Similarly, the node after $a_i$ is an exact incident edge $g_i$ colored
$a$ or $b$.  This first edge will certify below that the outside endpoint
is not isolated after the recoloring.  These conclusions use actual edge
nodes, so parallel edge identities are not identified.

By the terminal-kind hypothesis we are in the PE case.
The parallel field of \noderef{HierarchyBase} gives a support position,
the support vertex
\[
                 x=T_n.\operatorname{vertex}(\text{that position}),
\]
an exact prefix $T_n|x$, a tree $R_n$ which is a
\noderef{TashkinovClosure} of that prefix under $\psi$, and a witness
that $A$ is the \noderef{OrderedTreeUnion} $T_n\vee R_n$.  In
particular,
\[
 V(A)=V(T_n)\cup V(R_n),                              \tag{5}
\]
and $x\in V(T_n)$ by \noderef{TreeSequenceVertexSet}.

We verify separately that $x\in V(R_n)$.  By
\noderef{TreeSequencePrefix}, the prefix through the displayed position
contains its last vertex $x$.  Unfolding \noderef{TashkinovClosure}
gives a \noderef{TashkinovAugmentationSequence} from that prefix to
$R_n$.  Induct along this reflexive-transitive sequence.  In the
reflexive case the assertion is the prefix membership just noted.  In a
successor step, \noderef{TashkinovAugmentation} supplies a
\noderef{TreeSequenceInitialSegment} from the preceding tree to the next
one.  If the old occurrence of $x$ is at position $j$, the entrywise
vertex equality in that initial-segment relation identifies the new
vertex at the included position $j$ with $x$; then
\noderef{TreeSequenceVertexSet} puts $x$ in the next tree.  Induction
therefore gives
\[
                     x\in V(T_n)\cap V(R_n).          \tag{6}
\]

At this particular vertex $x$, let $C_x$ be the finite set
\noderef{PresentColors}$(G,\psi,x)$.  For every $d\in C_x$, choose the
incident exact edge supplied by that definition.  Properness makes an
edge of color $d$ at $x$ unique, and chosen edges for distinct colors are
distinct because $\psi$ is a function.  This injects $C_x$ into the
incident-edge set.  By \noderef{Degree} and \noderef{MaximumDegree},
\[
 |C_x|\leq\deg_G(x)\leq\Delta(G).
\]
The initial \noderef{KTriple} gives $\Delta(G)+1\leq k$, whereas the
palette $[k]$ has cardinality $k$.  Hence $C_x$ is a proper subset of the
palette.  Choose $c$ in its complement.  The definition
\noderef{MissingColors} gives $c\in\overline\psi(x)$, and (6) with
\noderef{TreeMissingColors} makes $c$ missing on both $T_n$ and $R_n$.

The pre-step vertex set $V(T_n)$ is elementary under $\psi_{n-1}$.  For
$n>1$, this is lower assertion (i) applied to the actual predecessor from
\noderef{CertifiedETTPredecessor}, whose tree and generating-coloring
equalities give exactly $T_n$ and $\psi_{n-1}$.  For $n=1$, it is the
elementarity conclusion of \noderef{TechnicalTheoremBase} for the
initial tuple.  All hypotheses of
\noderef{ParallelCommonMissingClosed} are now exact: the certified
terminal step is PE, (5) is the chosen union, and the pre-step
elementarity has just been supplied.  Since $c$ is common-missing, that
lemma gives
\[
                 \partial_{\psi,c}(A)=\varnothing.   \tag{7}
\]
Equivalently, \noderef{TreeColorSetClosed} says that $\{c\}$ is closed
on $A$.

The color $b$ is present at $x$.  If it were missing, then (6) would make
it common-missing on $T_n$ and $R_n$.  A second application of
\noderef{ParallelCommonMissingClosed} would give
$\partial_{\psi,b}(A)=\varnothing$.  The assumed singleton closedness of
$a$ gives $\partial_{\psi,a}(A)=\varnothing$, contradicting the exact
boundary identity (1).  Hence $b$ is not in
\noderef{MissingColors} at $x$, so it is in
\noderef{PresentColors} there.  Since $c$ is missing at the same vertex,
we again have
\[
                         c\ne b.                    \tag{8}
\]

We now define $\rho$, again separating equality from a genuine Kempe
change.  If $a=c$, put $\rho=\psi$.  The assumed properness and
one-uncoloredness persist.  Reflexivity in \noderef{sc2} makes $\rho$
stable on $T_n$ with control $D_n$ and on $R_n$ with empty control;
by definition these two clauses are exactly
\noderef{CombinedStableColoring}.  The two original walks, with $a$
renamed $c$, retain all their exit-path fields and disjoint supports.

Suppose $a\ne c$.  In the $(a,c)$ incidence graph define explicitly
\[
 \mathcal R_{\mathrm{PE}}=
 \{\operatorname{inl}a_1,\operatorname{inl}a_2\}
 \mathbin\cup
 \{\operatorname{inr}f: \operatorname{inr}f\in
       \operatorname{supp}P_1\cup\operatorname{supp}P_2,
       \ \psi(f)=a\}.                               \tag{9}
\]
This is a finite root set because the two walk supports are finite.  We
check that each rooted full component avoids $A$.  The endpoint roots
$a_i$ are outside $A$.
An $a$-colored edge root on $P_i$ cannot be the final edge because $a$ is
closed on $A$, so its component contains an earlier outside vertex of
the walk.  If a selected component also contained a vertex of $A$, a
shortest incidence path from such an outside vertex to the first vertex
of $A$ would have, between its last outside vertex and first inside
vertex, an exact boundary identity colored $a$ or $c$.  Unfolding
\noderef{TwoColorIncidenceGraph}, \noderef{BoundaryEdges}, and
\noderef{BoundaryColorEdges} gives an element of
$\partial_{\psi,a}(A)$ or $\partial_{\psi,c}(A)$, contrary respectively
to the given closedness or (7).  Thus the root set (9) satisfies the outside
condition in \noderef{OutsideKempeRecoloring} relative to $A$.

Swap $a,c$ on exactly these components to obtain $\rho$.  Equation (5)
shows $V(T_n),V(R_n)\subseteq V(A)$.  Therefore every selected component
also avoids each constituent, and the same roots and edge formula certify
\noderef{OutsideKempeRecoloring} relative separately to $A$, $T_n$, and
$R_n$.  Applying \noderef{OutsideKempeRecoloringStable} in these three
forms gives: $\rho$ is proper and leaves exactly $e_0$ uncolored; every
edge incident with $A$ is fixed; $\rho$ is stable on $T_n$ for the exact
control $D_n$; and it is stable on $R_n$ for the empty control.  The
last two clauses are precisely \noderef{CombinedStableColoring}.

For completeness, the transformed exit paths are checked directly, not by
preserving their old full components.  Every old $a$-edge node on either
walk was included in the root set (9) and
therefore has new color $c$; every old $b$-edge is fixed because (8) and
$a\ne b$ put $b$ outside the exchanged pair.  Consequently the same node
lists have all required adjacencies in $I_G(\rho;c,b)$.  Their simplicity,
length, endpoints, support disjointness, and unique intersections with
$A$ are unchanged.  If the old outside alternative was that $b$ is
missing, $b$ remains missing.  If it was that $a$ is missing, the root at
$a_i$ swaps every incident $a$- or $c$-edge in its full component, so $c$
is missing there under $\rho$.  The witness and complement definitions
\noderef{PresentColors} and \noderef{MissingColors} justify these last
two statements.  Thus all six fields in \noderef{TreeExitPath} hold for
both transformed walks.

Before using terminal PE interchangeability, we verify its two premises
under the changed coloring.  In the identity case, the previously chosen
$c$ is missing at $x\in V(A)$ and (7) is its empty colored boundary;
substitution of $\rho=\psi$ therefore gives both premises.  In the
genuine recoloring case, stability on
$T_n$ preserves the missing-color set at $x$ by the vertex clause of
\noderef{StableColoring}.  Hence $c$ remains missing at $x\in V(A)$, and
\noderef{TreeMissingColors} gives
$c\in\overline\rho(A)$.  For closedness, suppose
$f\in\partial_{\rho,c}(A)$.  Then \noderef{BoundaryColorEdges} gives
$f\in\partial_G(A)$ and $\rho(f)=c$.  Such a boundary identity is
incident with $A$, so the fixed-incident-edge conclusion of
\noderef{OutsideKempeRecoloringStable} gives
$\psi(f)=\rho(f)=c$.  The same underlying membership
$f\in\partial_G(A)$ then puts $f$ in
$\partial_{\psi,c}(A)$, contradicting (7).  Therefore
\[
 c\in\overline\rho(A)
 \quad\hbox{and}\quad
 \partial_{\rho,c}(A)=\varnothing.                  \tag{10}
\]
By \noderef{TreeColorSetClosed}, the second assertion is exactly
singleton $c$-closedness under $\rho$.

The three-leaf argument must likewise be made in the new incidence graph.
Combined stability at the common support $x$ preserves both $c$-missingness
and $b$-presence.  Properness therefore makes
$\operatorname{inl}x$ an exact one-neighbor node of
$I_G(\rho;c,b)$.  At each $a_i$, the first edge of the transformed exit
path supplies an incident edge in the nonmissing one of $c,b$, while its
missing-color alternative and properness exclude every other
$(c,b)$-edge.  Hence $\operatorname{inl}a_1$ and
$\operatorname{inl}a_2$ are exact one-neighbor nodes as well.

Their full \noderef{TwoColorComponent}s are path components by (8) and
\noderef{TwoColorPathComponent}.  The component through $x$ meets $A$ at
$x$, and the transformed walk from $a_i$ to $b_i$ proves that the
component through $a_i$ meets $A$ at $b_i$.  The three roots are distinct
because $a_1\ne a_2$ and the $a_i$ lie outside $A$ whereas $x$ lies
inside.  If all three full components were equal,
\noderef{TwoColorComponentSpanningPath} would give a simple spanning path
containing three distinct one-neighbor nodes.  Each would have to be one
of the two endpoints, which is impossible.  Thus the component through
$x$ differs from at least one outside-end component, and both are full
$(c,b)$ path components meeting $A$.

Finally apply conclusion (iii) of \noderef{interchange} to the actual
terminal PE, the given maximum property and lower-rung package, the exact
base and union, and the \noderef{CombinedStableColoring} just proved.
Its hypotheses are precisely (10), so it says that $c,b$ are
$A$-interchangeable under $\rho$.  Unfolding
\noderef{TreeInterchangeable} forces the two distinct path components
meeting $A$ to be equal, a contradiction.  This proves the PE branch
without ever constructing a zero-rung ETT.
\end{proof}
